\cotefile{149}

\scanpage{54}
par $E_{\bar K,K}$) en prenant son centralisateur $Z(L_i)$ dans $E$, d'où
\diagram{\begin{tikzcd}[ampersand replacement=\&,column sep=small]
(2)\qquad 1\arrow[r,no head]\&L_i\arrow[r,no head]\arrow[d,phantom,"\simeq"{description,rotate=90}]\&Z(L_i)\arrow[r]\&\Gamma\arrow[r,phantom,"{\unsure{\longrightarrow|-}}"{description}]\&1\\
\&T_\infty(\bar K^*)=T\&\&\&
\end{tikzcd}}
et on prend les scindages de cette extension. Il en existe, p.ex. définis par une uniformisante de $\underline{\mathcal O}_{\widehat U,i}$. L'ens des classes de conjugaison des \deleted{\unsure{ses|uns}} scindages \deleted{\unsure{les|de}} de (\unsure{2|1}) est un torseur sous \deleted{$H^1(\Gamma,T)$}
\[
(3)\qquad H^1(\Gamma,T)\simeq\varprojlim_n H^1(K,{}_n\mu)\simeq\widehat{K^*}=\varprojlim_n K^*_n
\]
et \deleted{\unsure{cet ensemble}} s'envoie injectivement \deleted{\unsure{vers}} dans l'ens des classes de $\pi$-conj. de scindages.

\uline{Proposition} On suppose $(g,\nu)\ne(0,0),(01)$ i.e. $\Pi_{\bar U}=\Pi_{U\unsure{/|\mid}K}\ne(1)$. Alors les classes de \inserted{$\pi$-conjugaison} scindages \deleted{\unsure{assoc.}} de (1) définis par les scindages de (2) sont distinctes de celles associées aux pts de $U(K)$. \deleted{De plus}, Si \inserted{de plus\unsure{,}} $(g,\nu)\ne(0,2)$, \deleted{\unsure{dans le}} \inserted{i.e. si on est} dans le cas anabélien, alors les classes de $\pi$-conjugaison de scindages de (1), associés à des scindages de (2) pour deux indices $i\deleted{\unsure{a}}=i_1$ et $i=i_2$ distincts\unsure{,} sont distinctes.
\margin{C'est démontré sauf pour le type $(0,3)$, \unsure{qui|que} mérite un examen approfondi\unsure{.}}
\margin{53}

\scanpage{55}
\mspage{27}
La première assertion s'obtient en ``bouchant''
le trou $i$, alors la section envisagée devient
la section de $U\cup\{i\}=U'$ associée au pt $i$,
et elle est donc distincte de celle \unsure{associée|associé}
aux \deleted{\unsure{pts|pt}} autres pts de $U'$, i.e. aux pts de
$U$ --- a fortiori avant de passer au
quotient par le ss-groupe inv. engendré
par $L_i$. On \unsure{raisonne|voudra} de $=$ fa\unsure{ç|c}on \unsure{se|en}
convaincre que les classes de conj. \unsure{des|de}
scindages associés à un $L_{i_1}$ et un $L_{i_2}$, $i_1\ne i_2$,
sont distincts --- Ça marche dans tous les
cas sauf le cas du type $(0,3)$, ---
car en bouchant deux trous, on tombe
sur le type $(0,1)$, où on ne dispose pas
du résultat d'injectivité. Mais on
peut s'en tirer autrement, à condition
d'admettre que pour un scindage de
(2), faisant opérer $\Gamma$ sur $\pi$, on a
\[
(3)\qquad\pi^{\Gamma^\circ}=L_i
\]
(donc $\pi^\Gamma=(1)$, d'ailleurs) --- résultat
qui me paraît hautement plausible. Cela
\margin{\figure{\begin{tikzpicture}\draw (0,0)--(0,-1.9);\draw (0.12,-0.5)--(0.12,-1.8);\end{tikzpicture}}}
\margin{À énoncer parmi les conjectures-clefs\ldots}
\margin{54}

\scanpage{56}
montrerait que la classe de conjugaison des sections détermine la classe de conjugaison de $L_i$, donc $i$.

\uline{Conjecture A} Soit $U$ courbe alg. anabélienne géom. connexe sur corps $K$ de type fini sur $\mathbb{Q}$. Alors toute section de (1) est \deleted{\unsure{définie|definie}} d'un des deux types précédents, i.e. soit définie par un pt de $U(K)$, soit par une section d'une extension (2), avec $i\in I(\bar K)^\Gamma$ i.e $i$ un pt de $\widehat U\setminus U$, rationnel sur $K$.
\margin{Il y a \inserted{\unsure{de plus|bcp} \unsure{des|de}} plus d'\illegible{1} les précisions $\pi^{\Gamma^\circ}=(1)$ (premier cas) $\pi^{\Gamma^\circ}=L_i$ (deux. cas)}

Si on admet cette conjecture, alors on en conclurait, par passage à la limite, en considérant le corps des fonctions $L$ de $U$ et $E_L\longrightarrow \unsure{E|B}_K$ ($E_L$ peut \unsure{être|etre} considéré comme un groupe à lacets ``infini'' car une infinité de \inserted{classes de} ss-groupes lacets\inserted{ $L_i$}\ldots) que tout\deleted{\unsure{e}} scindage de cette extension provient d'un scindage d'une extension du type (2), avec $i\in I^\Gamma=\unsure{X}(K)$ ($X=\widehat U$). Les classes \inserted{\unsure{conj.}} de tels scindages
\margin{55}

\scanpage{57}
\mspage{28}
se grouperaient donc \inserted{\unsure{classes de|\deleted{\unsure{classes} de}}} \unsure{\deleted{\unsure{par}\ \unsure{des}}|par des} paquets
(en regardant les centralisateurs des ss-groupes
images de $\Gamma^\circ$ par ses sections,\deleted{\unsure{)|f|et}} et \unsure{en|un}
mettant ensemble des scindages qui ont des
$\unsure{\mathrm{Zentr}|\mathrm{Centr}}(\Gamma^\circ)$ conjugués (même s'ils \unsure{ne|en} sont
eux-mêmes conjugués). Donc on retrouve
encore une description de $X(K)$ (donc
aussi des $X(K')$ pour tt ext. \unsure{finie|fini} \inserted{$K'$} de $K$)
en termes de l'extension $E_L$ de $E_K$ par
$\Pi_{L/K}$, en $\unsure{=|\overline m}$ temps qu'une façon de
reconstituer les $U=X-\overline{\unsure{I|\underline I}}$\ldots
Donc en fait c'est la structure $E_L\unsure{-|\longrightarrow}E_K$
qui est la plus riche a priori, et de
loin plus commode pour le genre $0$ et $1$,
où la considération des $U$ de type $(g,\nu)$\unsure{.|,}
$(2g\unsure{+}\nu\geqslant3)$ casse le groupe ``continu''
d'automorphismes\ldots La forme ``modélique''
de la conjecture précédente revient
à la forme ``birationnelle'', quand on y
précise cette dernière en disant que tt
\margin{56}

\scanpage{58}
scindage de $E_U-E_K$ se \deleted{\illegible{1}} remonte en un scindage de $E_L-E_K$ (ou aussi, \deleted{\unsure{par|pour}} en un scindage de $E_V-E_K$, si $V$ est un modèle $\subset U$)

On va mettre les conjectures précédentes (sous forme modélique, disons) sous une forme plus géométrique, en introduisant, en $\unsure{=|\overline m}$ temps qu'un revêtement universel \deleted{\unsure{de}} $\widetilde U$ de $U$, le normalisé \unsure{\deleted{de}|\inserted{\deleted{de}}|\inserted{de}} $X'$ de $X$ dans $\widetilde U$ (où $X=\widehat U$). (NB je m'abstiens de le noter $\widetilde X$, vu que justement il n'est pas étale sur $X$\ldots). Notons que pour $i\in\bar I=\bar X-\bar U$, les \inserted{l'\unsure{ens|un} des} $L_i$ de $\bar\pi=\pi(\bar U)$ sont en corr. $1-1$ avec l'e\deleted{\unsure{ns}}\unsure{.|\deleted{.}} fibre $X'_i$ de $X'$ au dessus de $i$

\diagram{\begin{tikzcd}[ampersand replacement=\&,column sep=large,row sep=large]
X\arrow[r,no head]\arrow[d,phantom,"\cup" description]\&\bar X\arrow[r,no head]\arrow[d,phantom,"\cup" description]\&X'\arrow[d,phantom,"\cup" description]\\
U\arrow[r,no head]\arrow[d,no head]\&\bar U\arrow[r,no head]\arrow[d,no head]\&\widetilde U\\
K\arrow[r,no head]\&\bar K\&
\end{tikzcd}}

Donc $X'$ peut être considéré comme la réunion de $\widetilde U$, et de $X'|\bar I=$ ensemble des sous-groupes lacets de $\bar\pi$, qui apparaissent ainsi comme des ``pts \unsure{à|:} l'infini'' du rev. universel $\widetilde U$. D'ailleurs $E_U$ s'inter-
\margin{57}

\scanpage{59}
\mspage{29}
prète comme le groupe \inserted{\unsure{compatibles} avec} des $\unsure{U|\deleted{K}|\deleted{U}}$-autom du schéma $\widetilde U$, et $E_U\longrightarrow E_K$ comme l'hom de passage au quotient évident (\uline{NB} $\bar K$ s'identifie à la clôture alg. de $K$ dans $H^0(\widetilde U,\underline{\mathcal O}_{\widetilde U})$, donc $E_U$ opère sur $\operatorname{Spec}\bar K$ de façon naturelle\ldots) Une section de $E_U$ sur $E_K$ est donc une action de $E_K$ sur $\widetilde U$, compatible avec son action sur $\widetilde U$, et continue dans un sens convenable (sans doute la conv. simple-discr\unsure{è|e}te sur les \deleted{$\unsure{\bar U_i^{a}|\widetilde U_i^{a}|\bar U_i^{\alpha}|\widetilde U_i}$} $\bar U_i$ finis sur $\bar U$ entre $\widetilde U$ et $\bar U$\ldots).

Considérons donc la

\uline{Conjecture} \uline{B} Toute telle action de $\Gamma$ sur $\widetilde U$ admet \uline{dans $X'$}$=\widehat{\widetilde U}$ un pt fixe et un seul.
\margin{\figure{\begin{tikzpicture}\draw (0,0)--(0,-3.1);\end{tikzpicture}}}
\margin{et $\unsure{=|\overline m}$ l'action induite de $\Gamma^\circ$ doit avoir un pt fixe et un seul --- cf plus bas.}

Ceci signifie donc

a) S'il y a un pt fixe à distance finie ie. $\widetilde x\in\widetilde U^\Gamma$, alors 1°) L'image de $\widetilde x$ dans $U$ est uniquement déterminée --- c'est essentiellement le th.3 du $\S1$ (deux pts distincts de $U(K)$ définissent des classes de conj. de scindages distinct\unsure{es|s}) et
\margin{58}

\scanpage{60}
2°) $\pi^\Gamma=(1)$ (i\unsure{.|:}e il n'y a pas d'autres
p\textsuperscript{t} fixes dans $\widetilde U$ sur le $\unsure{=|\overline m}$ $x\in U(K)$,
et enfin

3°) il n'y a pas en $\unsure{=|\overline m}$ temps de pt
fixe $\deleted{\unsure{x|\mathrm{en}}}\ \widetilde x$ à l'infini --- i.e il n'existe
pas de $L_i$ normalisé par $\Gamma$, i.e. un scindage
du premier type n'est pas en $\unsure{=|\overline m}$ temps
du deuxième (fait que nous avons
établi directement, précédemment).

\margin{C'est un cas particulier d'une ``conséquence conjecturale'' des conjectures du $\S1$ --- sauf qu'il faudrait $\unsure{=|\overline m}$ écrire $\pi^{\Gamma^\circ}=(1)$}

D'autre part, dans le cas de pts fixes
à l'$\infty$, l'unicité de l'image dans $X$
signifie qu'une $\unsure{=|\overline m}$ action effective ne
peut normaliser à la fois un $L_i$ et un $L_j$
($i\ne j$) --- fait également établi sauf
dans le cas $(g,\nu)=(0,3)$ --- et l'unicité
--- dessus d'un $i\in\bar I$ fixé signifie
que le $L_i$ \inserted{($i$ fixé)} normalisé par $\Gamma$ est unique,
ce qui est un affaiblissement de la
relation $L_i=\unsure{\deleted{\unsure{Z}}\mathrm{er}|\deleted{\unsure{Z}}\mathrm{en}|\deleted{\unsure{Z}}_{\mathrm{en}}|\deleted{\unsure{Z}}_{\pi}}\pi^{\Gamma^\circ}$
pour ces opérations, conjecturée plus haut.
\margin{59}

\scanpage{61}
\mspage{30}
En fait, je conjecture que dans la conj. B, il est $\unsure{=|\overline m}$ vrai que $\Gamma^\circ$ agissant sur $X'=\widehat{\widetilde U}$ a un pt fixe et un seul (ce qui est plus fort, que ce point \inserted{fixe} sera nécessairement fixe par $\Gamma$). Ceci \deleted{\unsure{signifie}}\inserted{implique}, dans le cas \unsure{des|de} pts fixes à distance finie, qu'on a alors $\pi^{\Gamma^\circ}=(1)$, comme il se devrait en général --- et dans le cas de pt fixe à l'infini, que
\[
\pi^{\Gamma^\circ}\subset\operatorname{Norm}_{\pi}(L_i)=L_i
\]
donc la relation $\pi^{\Gamma_0}=L_i$ ``hautement plausible'' de tantôt !

Du coup je me demande d'ailleurs si ces énoncés ne devraient pas être valables sans partir d'une action de $\Gamma$ \deleted{\unsure{tt}} entier, mais seulement de $\Gamma^\circ$ (ou d'un sous-groupe contenant un sous-groupe ouvert de $\Gamma^\circ$, plus gén.) --- ce qui signifierait plus ou moins que
\margin{60}

\scanpage{62}
tous nos beaux énoncés devraient être valables, non seulement pour un corps de base $K$ ext. de type fini de $\mathbb{Q}$, mais dès que $K$ est ext. de type fini d'un corps cyclotomique (pas néc. fini sur $\mathbb{Q}$).

Nous pourrions définir les \uline{courbes} de \uline{Poincaré} sur un corps alg. clos $\bar K$ de car $0$, comme étant les courbes isomorphes à des rev. universels de courbes alg. anabéliennes sur $\bar K$ (deux courbes anabéliennes $=\ \bar U,\bar V$ sur $\bar K$ définiss\unsure{e|a}nt des rev. de Poincaré isomorphes, ssi $\exists$ un revêtement fini étale de l'une, isomorphe à un rev. fini étale de l'autre\ldots).

\unsure{E|É}tant donné une courbe de Poincaré\unsure{.} $\unsure{\widetilde{\mathcal U}|\widetilde{\overline{\mathcal U}}}$ sur $\bar K$, \deleted{telle} on définit canoniquement sa complétée\inserted{ $\widehat U$}, de telle façon que $\forall$ autom de $(\widetilde{\mathcal U},\bar K)$ (pas néc. induisant l'identité sur $\bar K$) induise un autom. de $\widehat{\mathcal U}$.
Ceci posé :
\margin{61}

\scanpage{63}
\mspage{31}
\uline{Conjecture B'} \unsure{Soient|Soit} $K$ un corps de type fini sur $\mathbb{Q}$ (ou sur un corps cyclotomique, cela suffira peut-être), $\bar K$ une clôture algébrique de $K$, $\mathcal U$ une courbe de Poincaré sur $\bar K$, de complétée $\widehat{\mathcal U}=\mathcal X$. Considérons une action de $\Gamma=\Gamma_{\bar K|K}$ sur $\mathcal U$, compatible avec son action sur $\bar K$, d'où une action de $\Gamma$ sur $\mathcal X$. Ceci posé :

il existe un point fixe et un seul de $\Gamma^\circ$ agissant sur $\mathcal X$

($\Gamma^\circ$, noyau du caractère cyclotomique $\Gamma\xrightarrow{\chi}\widehat{\mathbb Z}^{\!*}$).

La différence avec la conjecture B, que celle-ci renforce, c'est qu'ici on ne suppose \uline{pas} que l'action de $\Gamma$ sur $\mathcal U$ normalise l'action d'un groupe profini $\pi$, agissant librement sur $\mathcal U$ de façon \unsure{$\mathcal V$|que|q.} $\mathcal U/\pi$ soit une courbe alg. anabélienne sur $\bar K$.
\margin{62}

\scanpage{64}
Que deviendraient les conjectures précédentes, quand on les applique \unsure{à|:} une situation où $K$ est remplacé par un modèle $S$ de $K$ (disons, élémentaire anabélien), quand $U_K$ provient d'une courbe relative\unsure{.} $U_S$ sur $S$ — de sorte qu'on a un hom de groupes
\diagram{\begin{tikzcd}[ampersand replacement=\&,row sep=1.2em,column sep=6em]
{} \& E_{U_S}\arrow[dd] \\
(4) \& {} \\
{} \& E_S
\end{tikzcd}}
de noyau $\pi_{\bar U}$, dont $E_{U_K}-E_K$ est déduit par chgt de base i.e. par produit fibré
\diagram{\begin{tikzcd}[ampersand replacement=\&,row sep=1.2em,column sep=4em]
{} \& \pi\arrow[r,no head]\arrow[d] \& \pi\arrow[d] \\
{} \& E_{U_S}\arrow[dd] \& E_{U_K}\arrow[l]\arrow[dd] \\
(5) \& {} \& {} \\
{} \& E_S \& E_K\arrow[l,"\mathrm{surj}"']\arrow[uul,dashed,"\varphi"']
\end{tikzcd}}
Ainsi, les sections de $E_{U_{\deleted{\unsure{S}}K}}$ sur $E_{\deleted{\unsure{S}}K}$ correspondent aux \deleted{\uline{\unsure{fausses|bonnes}} sections de} \inserted{relèvements continus} $E_K\longrightarrow E_{U_S}$ de l'hom surjectif $E_K\longrightarrow E_S$, et parmi ces relèvements, ceux qui sont \deleted{\unsure{nuls}} \inserted{triviaux} sur le noyau de $E_K\longrightarrow E_S$ correspondent exactement aux sections
\margin{63}

\scanpage{66}
dans quelle mesure la conjecture A
(ou B) faite dans \unsure{le|ce} §, permet de
reconstruire la catégorie des modèles élém.
anabéliens sur $Q$, \inserted{et} celle des extensions
de type fini de $Q$ et des modèles
élém. anabéliens sur ceux-ci, en
termes des groupes extérieurs \inserted{(ou topos \unsure{galoisiens})} associés
à partir bien sûr de la donnée
fondamentale de $\Gamma_Q=\Gamma_{\bar Q/\bar Q_0}$, opérant
extérieurement sur $\widehat\pi_{03}$, d'où déjà
l'extension $E_{03}=E_{U_{03}/Q}$, où
$U_{0,3}=\mathbb P^1_Q-\{0,1,\infty\}$.

Prenons les \inserted{classes de $\widehat\pi_{03}$-conj. de \unsure{scissions|sections}} sections de $E_{03}$ sur $\Gamma_Q=\Gamma$
\inserted{telles que le centralisateur de $\Gamma^\circ$ soit trivial (sections ``admissibles)''},
i.e. les ``points'' \deleted{\unsure{de|des}} \inserted{(ou les sections) du} topos $B_{\widehat\pi_{0,3},\Gamma_Q}^{\natural}$ sur
$B_{\Gamma_Q}^{\natural}$ — on trouve un ensemble sur
lequel $\Gamma$ opère (qui n'est autre que
$U_{03}(\bar Q_0)$, à deux cas. près). Pour
tout ens. fini. \inserted{$I$} de sections, \inserted{stable par $\Gamma$,} le foncteur
\margin{65}

\scanpage{67}
\mspage{33}
``forage de trous'' doit nous fournir un groupe extérieur $\Pi_{0,3}(I)$, \inserted{de type $(0,3+\operatorname{card} I)$,} sur lequel $\Gamma$ opère (il vaut mieux peut-être utiliser le yoga introduit par ailleurs des groupoïdes rigides — donc on \unsure{peut|pourra} ou devra aussi boucher certains des trous $0,1,\infty$ — on trouve donc l'équivalent groupoïdal de la \deleted{``sphère \unsure{ar.}} droite projective $\mathbb P^1_Q$, on l'appelle la sphère arithmétique — qui correspond à un groupe \inserted{ext} à lacets infini sur lequel $\Gamma$ opère — en fait, ce n'est autre que $E_{K_1}$, où $\vert$
\[
\tag{6}K_1=Q(\deleted{\unsure{X_1|x_1}}\,T_1)
\]
est l'extension transcendante pure \inserted{typ.} de degré $1$ de $Q$.

Partant de (6), on construit de $\overline m$ l'équivalent groupoïdal de $\deleted{\unsure{X|P}}\,U_{0,3}$ \unsure{$\longleftarrow$|\deleted{(\unsure{qu}}} et on remonte comme
\margin{66}

\scanpage{68}
précédemment, pour avoir, soit des courbes de type $(0,\nu_2)$ sur $K_1$\unsure{,} (ou sur une extension finie de $K_1$) soit des courbes relatives de type $(0,\nu_2)$ sur une courbe sur $Q$ \deleted{de type} (ou une extension finie de $Q$, ou un rev. \deleted{étale} étale fini d'une telle $U_{0\nu_1}$).

On procède de proche en proche pour construire finalt tous les $E_K$ sur $\deleted{\unsure{E}}\,E_Q$ (car tt corps ext. de type fini de $Q$, est ext finie d'une ext. transcendante pure) et tous les modèles élémentaires, où à chaque cran la fibration est soit une courbe de genre $0$, soit un rev. étale fini d'une telle fibration particulière. Sauf erreur, ça fait assez pour avoir un syst. fondamental de voisinages de tt point \unsure{d'une|d'un} $X$ lisse sur un $K$
\margin{67}

\scanpage{70}
Reprenons le cas de $U=U_S$ schém. relatif sur $S$, ``élémentaire'' sur $S$ -- à fibres successives anabéliennes (s'il le faut) ou du moins à $\pi_1$ non nul, $S$ étant lui-$\overline m$ (pour fixer les idées) lisse sur $\mathbb{Q}$, irréductible, corps de fonctions $K$, et reprenons le diagramme 5. Considérons une section rationnelle \inserted{$f$} de $U$ sur $S$, définissant une section de $E_{U_{\unsure{\eta|K}}}$ sur $E_K$ -- ou, ce qui revient au même, un relèvement de l'hom surjectif $E_K\twoheadrightarrow E_S$ en $E_K\longrightarrow E_{U_S}$ \inserted{(composé de \inserted{la section} $E_K-E_{U_{\unsure{\eta|K}}}$ avec $E_{U_K}\unsure{-|\longrightarrow}E_{U_S}$)}. Je veux montrer que $f$ est partout \unsure{défini|defini} i.e. une section \deleted{\unsure{rationnelle}} de $U_S$ sur $S$, ssi la section de $E_{U_{\unsure{\eta|K}}}$ sur $E_K$ provient d'une section de $E_{U_S}$ sur $E_S$, i.e. ssi le relèvement $E_K\longrightarrow E_{U_S}$ s'annule sur
\margin{69}

\scanpage{71}
\mspage{35}
le noyau de $E_K\longrightarrow E_S$.

Notons que cette dernière condition est une condition ``de codimension 1 sur $S$'' -- de façon plus précise, si $Z$ est un sous-schéma fermé de $S$ de codim $\geq 2$, \deleted{alors} alors, posant $S'=S\setminus Z$, on a $\pi_1(S')\xrightarrow{\sim}\pi_1(S)=\unsure{E|\deleted{F}E}_S$ par le ``théorème de pureté'' -- donc le noyau de $\deleted{\pi_1}E_K\longrightarrow E_S$ est le $\unsure{=|\overline{m}}$ que celui de $E_K-E_{S'}$, ou, si on préfère (comme $S'$ n'est pas un ``modèle'') que le sous-groupe \deleted{\unsure{nor-}}\inserted{fermé} engendré par les noyaux des $E_K\longrightarrow E_{S'_i}$, où les $S'_i$ sont des ouverts ``modèles'' qui recouvrent $S'$. Si donc la condition envisagée est remplie relativement aux $S'_i$ (qui pour-\deleted{\unsure{tant}} ne recouvrent pas $S$, seulement $S'$) -- ce qui est censé \unsure{signifier|signifie} que la section
\margin{70}

\scanpage{72}
rationnelle envisagée est morphique sur
les $S'_i$, i.e. sur $S'$ -- alors elle est
vérifiée relativement \unsure{à|:} $S$ -- ce qui
est censé signifier que la section est
morphique sur $S$. Donc, pour se
rassurer, il faudrait s'assurer a
priori qu'une section de $U_{S'}$ sur $S'$
se prolonge automatiquement en \deleted{$U$}
une section de $U_S$ sur $S$. De proche
en proche, on est ramené au cas
d'une courbe relative $U_S=X_S-T$,
$X$ lisse \inserted{sur $S$} de dim. rel\unsure{a|.}tive 1, $T$ fini étale
sur $S$, \deleted{\unsure{le}} OPS $\overline m$ $T$ décomposé sur $S$.
Si $X$ de genre relatif $\geq 1$, on sait
(th. de Weil) que la section se prolonge
en une section de $X$, soit $D$ l'image
inverse de $T$, c'est un diviseur sur $S$,
dont la trace sur $S'=S\setminus Z$ est nul, donc
(comme $\operatorname{codim}(Z,S)\geq 1$) il est nul, OK.
\margin{71}

\scanpage{73}
\mspage{36}
\diagram{\((9)\qquad\begin{tikzcd}[ampersand replacement=\&,column sep=small]
U_S \arrow[r,no head] \arrow[d,no head] \& U_{O_x} \arrow[r,no head] \arrow[d,no head] \& U_K \arrow[r,no head] \arrow[d,no head] \& U_{\overline K} \arrow[r,no head] \arrow[d,no head] \& \widetilde U_S \\
S \arrow[r,no head] \& \underline O_x \arrow[r,no head] \& K \arrow[r,no head] \& \overline K \&
\end{tikzcd}\)}
\diagram{\((10)\qquad\begin{tikzcd}[ampersand replacement=\&]
E_{U_S} \arrow[d] \& E_{U_{O_x}} \arrow[l] \arrow[d] \& E_{U_K} \arrow[l] \arrow[d] \\
E_S \& E_{\underline O_x} \arrow[l] \& E_K \arrow[l] \arrow[ul,dashed,"\varphi_x"']
\end{tikzcd}\)}
avec des carrés cartésiens, et des flèches horizontales surjectives. L'hom $E_K\longleftarrow E_{U_S}$ est composé d'un relèvement $E_K-E_{U_{O_x}}$ de $E_K\longrightarrow E_{O_x}$, avec l'hom canonique $\widehat E_{U_{O_x}}-E_{U_S}$. (Les \inserted{\deleted{\unsure{derniers|dernier}}} relèvements $E_K\longrightarrow E_{U_{O_x}}$ correspondent biuniv.$^{\mathrm t}$ aux sections de $E_{U_{\unsure{\eta|K}}}$ sur $E_K$, ou aux rel$^{\mathrm t}_{\mathrm s}$ de $E_K-E_S$ en $E_K-E_{U_S}$\ldots)

Ceci dit, j'ai envie de prouver que $\varphi_x:E_K\longrightarrow E_{U_{O_x}}$ s'annule \unsure{sur|dans} $I_x$, i.e. provient d'une section de $E_{U_{O_x}}$ sur $\underline O_x$ sss la section rationnelle corresp.$^{\mathrm t}$ * de $U_S/S$ est définie en $x$. Ceci impliquera
\margin{\uline{NB} Le quotient de $E_K$ par le \unsure{s.}groupe engendré par $\unsure{I_x|\overline I_x}$ n'est autre que $E_{\underline O_x}=\pi_1(\operatorname{Spec}\unsure{O_x|\underline O_x})$}
\margin{72}

\scanpage{74}
l'assertion précédente (que la section $\psi$ de $E_{U_\eta}$ sur $E_K$ provient d'une section de $E_{U_S}$ sur $E_S$, sss la section rationnelle correspondante est morphique)

Mais il s'agit ici d'un énoncé en fait purement géométrique, que j'ai envie de reformuler sous forme plus générale :

\uline{Théorème.} \deleted{\unsure{2|3}}\unsure{(?)} Soit $\unsure{\deleted{S}T|\deleted{T}S}$ un trait (spectre d'un anneau-jauge), $U$ un schéma relatif ``élémentaire'' sur $T$, anabélien*, $K$ le corps des fonctions de $T$, On choisit un rev. universel $\widetilde U$ de $U$, d'où une clôture alg. $\overline K$ de $K$, et on considère l'extension
\[
E_U=\pi_1(U;\widetilde U)\quad\text{de}\quad E_K=\pi_1(K,\overline K)\simeq\operatorname{Gal}(\overline K/K)
\]
par $\pi=\pi_1(U_{\overline K},\widetilde U)$. \deleted{Considérons une} \inserted{On a donc} \deleted{section pt $K$-rationnel de $U_K$}\unsure{\deleted{,}|,} \inserted{un carré cartésien de groupes profinis}

\uline{d'où}
\diagram{\begin{tikzcd}[ampersand replacement=\&]
E_U \arrow[d,"\mathrm{surj}"'] \& E_{U_\eta} \arrow[l,"\mathrm{surj}"'] \arrow[d,"\mathrm{surj}"] \\
E_S \& E_K=\operatorname{Gal}(\overline K/K) \arrow[l,"\mathrm{surj}"]
\end{tikzcd}}
où $E_S$ s'identifie au quotient de $E_K$ par le sous-groupe nor. fermé engendré par un groupe
\margin{* anabélien est inutile -- il suffit que les fibres de dim 1 des ``fibrations élémentaires'' de $U$ ne soient pas des \unsure{typ|types} $(0,0)$ ou $(0,1)$ -- i.e. à $\pi_1$ nul}
\margin{73}

\scanpage{75}
\mspage{37}
d'inertie $I_{x\unsure{'}}\unsure{\simeq|=}T_\infty(\overline K)$, cf plus haut. Soit $f_K$ $K$ un pt de $U_K$ rat/$K$, d'où une section $\psi=\psi_{f_K}$ de $E_K$ sur $E_{U_K}$. Ceci posé les conditions suivantes sont équivalentes

\noindent a) $f_K$ se prolonge en une section de $U$ sur $S$

\noindent b) $\psi$ provient d'une section de $E_U$ sur $E_S$

\noindent b') le composé $E_K\xrightarrow{\psi}E_{U_K}\unsure{\longrightarrow|-}E_U$ s'annule sur $I_x$

\deleted{\unsure{c)}}

L'équivalence de b) et b'), et qu'ils soient impliques par a), est clair. C'est l'implication $\mathrm{b)}\Longrightarrow\mathrm{a)}$ qui demande une démonstration. On est ramené aussitôt, par descente, au cas \deleted{\unsure{d'un}} où $T$ est strictement local (donc $E_K=\operatorname{Gal}\overline K\unsure{\mid|/}K$ est réduit à son ss-groupe d'inertie, \inserted{et $E_S=(1)$}). On est ramené de proche en proche au cas où $U/S$ est une courbe relative élémentaire, \deleted{de} \inserted{\unsure{le|en}} \deleted{genre $g\geq 1$}\unsure{\deleted{,}|,} \inserted{$U=X\setminus T$, $X$ lisse et propre.} Alors $f$ se prolonge en une section $f$ \deleted{rationnelle} \inserted{de} $X$ sur $S$, et la conclusion est que $f(S)\subset U$. Donc on est ramené finale-\newline ment au

\uline{Lemme} Soit $X$ schéma projectif lisse de dim relative 1 connexe sur $\unsure{\deleted{T}S|\deleted{S}T}$ trait strictement local,
\margin{74}

\scanpage{76}
Soit $T\subset X$ sous schéma, fini étale sur $S$, donc $T\simeq I_S$, $I$ ensemble fini, et soit $U=X\setminus T$ (donc $T$ est défini par une famille $(g_i)_{i\in I}$ de sections disjointes de $X$ sur $S$). \inserted{Si $g$ est le genre relatif, $\nu=\operatorname{card}I$, on suppose $(g,\nu)\neq(0,1)$} Soit \deleted{$\unsure{f_0}$} $i_0\in I$, $f$ une section \deleted{disjointe}\inserted{de $X/S$ distincte} des $g_i$, et telle que \deleted{$f\vert K(s)=g_{i_0}$} $f$ et $g_{i_0}$ coïncident en $s$ (pt fermé de $S$). \deleted{Considérons} \inserted{\deleted{Donc}} \inserted{Si $\eta=S\setminus\{s\}$,} \deleted{et le groupe de} \deleted{par l'homomorphisme} on a donc un morphisme $\eta\longrightarrow U$, d'où $\pi_1(\eta)\longrightarrow\pi_1(U)$. Je dis que cet hom. n'est pas trivial, \deleted{\unsure{donc}} et $\overline m$, \inserted{si $\nu\geq2$,} que pour \inserted{$\ell$ donné}
\diagram{\begin{tikzcd}[ampersand replacement=\&,column sep=small]
\& H_1(\eta,\mathbb Z_\ell) \arrow[r] \& H_1(U,\mathbb Z_\ell) \\
\inserted{T_\ell(\unsure{\overline K^*|\widehat K^*})} \arrow[ur,"\sim"] \& \& \inserted{[H_1(U_s,\mathbb Z_\ell)]} \arrow[u,"\sim"']
\end{tikzcd}}
n'est pas trivial (pour, \deleted{tout} $\ell=$ nb premier distinct de la caracteris-\newline
tique résiduelle).

Comme la section rationnelle de $\mathcal J^1_{X/S}$ définie par $f$ est régulière, on voit que le composé de l'hom envisagé avec $H_1(\mathcal J^1_{X/S},\mathbb Z_\ell)$ est nul -- i.e. le $H_1(\eta,\mathbb Z_\ell)$ s'envoie dans la partie torique de $H_1(U,\mathbb Z_\ell)$ $\bigl(\xleftarrow{\sim}H_1(U_s,\mathbb Z_\ell)\bigr)$, qui est canoniquement isomorphe à $\deleted{\mathbb Z}\inserted{T}_\ell^I/\deleted{\mathbb Z}\inserted{T}_\ell$.

[NB cette partie est nulle si $\operatorname{card}I=1$, et
\margin{75}

\scanpage{77}
\mspage{38}
dans ce cas le critère homologique s'avère insuffisant\ldots) Il faudrait donc calculer cet \deleted{élément de}\inserted{homomorphisme}
\[
(11)\qquad T_\ell\bigl(\simeq H_1(\eta,\mathbb Z_\ell)\bigr)\longrightarrow\unsure{T|\mathbb T}_\ell^I/T_\ell
\]
pour constater qu'il n'est pas nul dans le cas \unsure{envisagé|envisage}, $\nu\geqslant2$ (et traiter séparé-\newline
ment le cas $\nu=1$).

\deleted{[On va faire des calculs intrinsèques, en}
\[
\deleted{\unsure{\operatorname{Ker}}(\mathcal J^\bullet_{U/S}\longrightarrow\mathcal J^\bullet_{X/S})\simeq G_m^I/G_m}
\]
\[
\deleted{H_1(\mathcal U,\mathbb Z_\ell)\simeq H_1(\mathcal J^1_{U/S},\mathbb Z_\ell)\underset{\substack{\text{isom}\\\text{de Kummer}}}{\simeq}\lim_{\longleftarrow}\deleted{\unsure{\ell^n}}{}_{\ell^n}(\mathcal J^0_{U/S})}
\]
\deleted{On considère donc l'image inverse par $\eta\longrightarrow\mathcal J^1_{U/S}$ composé de $\eta\longrightarrow U\longrightarrow\mathcal J^1_{U/S}$) des isogénies}

Je vais deviner le résultat :

Soit \deleted{\unsure{q}} $x=g_{i_0}(s)$, $A=\underline O_{X,x}$, $V$ l'anneau de $S$, $\mathcal J_{i_0}$ l'idéal de l'hom $A\xrightarrow{g_{i_0}^*}V$ associé à $g_{i_0}$, c'est donc un idéal inversible de $A$ -- soit de $\overline m$ $\mathcal J_f$ l'idéal associé à $f^*:A\longrightarrow V$, et considérons $g_{i_0}^*(\mathcal J_f)$, c'est un idéal de $V$ engendré par un générateur, et comme $g_{i_0}\neq f$, on voit que cet idéal n'est pas nul. Soit $\mu=\operatorname{long}V/g_{i_0}^*(\mathcal J_f)\ \deleted{\unsure{=}}$. , cet entier
\margin{76}

\scanpage{78}
ne change pas quand on intervertit les roles de $g_{i_0}$ et $f$, car
\[
V/g_{i_0}^*(\mathcal J_f)\simeq A/(\mathcal J_{g_{i_0}}+\mathcal J_f)
\]
-- car $\nu$ est la longueur du schéma artinien local, intersection des \inserted{ss-schémas} $f(S)$ et $g_{i_0}(S)$ de $X$ -- on peut aussi l'interpréter comme une multiplicité d'intersection. Ceci posé, Je présume que l'hom
\[
T_\ell\longrightarrow T_\ell^I/T_\ell
\]
est le produit par $\mu$ de l'injection canonique $T_\ell-T_\ell^I$, correspondant à l'indice $i_0$. Il faudrait que je demande à Carlos de me le confirmer.

Reste le cas $\nu=1$, qui semble demander un traitement séparé.

\unsure{Le|\deleted{\illegible{1}}Le} \deleted{\illegible{1}}fait à vérifier (pour les groupes fondamentaux premi\unsure{e\uline{rs}|\uline{ers}}\margin{\deleted{\unsure{?!}}\\?!} à $p$) c'est que \unsure{le|l'} hom. ext. $\pi_1(\eta)\longrightarrow\pi_1(U)\simeq\pi_1(U_s)$ est égal à $\deleted{\unsure{\mu}}\unsure{k|\kappa}_{i_0}\circ(\mu.\operatorname{id}_T)$, où $\unsure{k|\kappa}_{i_0}$ est
\margin{Ceci doit être indépendant de la valeur de $\nu$ !}
\margin{77}
